\paragraph{\S105. Duality of presheaf topoi\newline and tentative additive generalizations (pp.\ 394--406).}
Dated 12 August with a handwritten addition of 13 August, subsection D explicitly limits its heading to topoi of the type $A^\wedge$ and tentative generalizations. Proposition~3 states the free-cocompletion property of the Yoneda inclusion $\varepsilon_A:A\to A^\wedge$: restriction gives (33), an equivalence between cocontinuous functors $A^\wedge\to\mathcal M$ and arbitrary functors $A\to\mathcal M$, when $\mathcal M$ has small inductive limits. The quasi-inverse sends $i$ to $i_!$, computed in (35) by the inductive limit of $i(a)$ over $A_{/F}$, and (36) gives its right adjoint $i^*(x)=(a\mapsto\mathrm{Hom}(i(a),x))$. The following remark literally assumes a pseudocotopos, where the witness says pseudotopos. Formula (37), however, writes $i\mapsto i^*$ as a covariant functor between the ordinary Hom categories without an opposite. Natural transformations reverse under this assignment, and the explicit opposite in (6), p.\ 385, rules out treating the missing variance mark as an unstated convention; the formula is preserved literally rather than mathematically amended (p.\ 394). The universal property adjoins arbitrary inductive limits while disregarding any that may already exist in $A$. Its projective-limit dual is (38), whose restriction map visibly includes $F\circ\varepsilon_A^o$, with the F and composition sign omitted by the witness. Thus $\mathcal M$-valued cosheaves on $A^\wedge$ identify with functors $A\to\mathcal M$, while $\mathcal N$-valued sheaves identify with functors $A^o\to\mathcal N$, equivalently presheaves on the chaotic site. If $\mathcal M$ has both kinds of small limits, combining these descriptions for $B=A^o$ gives (39), identifying cosheaves on $A^\wedge$ with sheaves on $B^\wedge$ with the same values. The final word canonically on p.\ 395 is canceled; diagram (40) then explains the comparison through $\mathrm{Hom}(A,\mathcal M)=\mathrm{Hom}(B^o,\mathcal M)$ and a chosen quasi-inverse supplied by projective limits. Its dual, with $\mathcal N=\mathcal M^o$, is canonically quasi-inverse to the comparison for B, using the ordinary opposite-category identity (6). The source leaves the two inner Hom arguments at the end of p.\ 396 blank, whereas the witness supplies them (pp.\ 395--396). Case~1 takes sets as coefficients. Formula (41) and its symmetric counterpart identify each presheaf topos with the set-valued cosheaves on the other, through the separately cocontinuous pairing (42), denoted $F*G$ in (43). Two applications of Proposition~3 give (44), so the pairing is determined on representables by $a*b^o=\mathrm{Hom}_A(b,a)=\mathrm{Hom}_B(a^o,b^o)$ in (45). Its symmetry compares the pairing for $(A,B)$ with the pairing for $(B,A)$; it is not an undeclared symmetry inside one fixed pair of categories. Formula (47) gives the evaluation identities $a*G\simeq G(a)$ and $F*b^o\simeq F(b)$ (pp.\ 396--398). The remarks postpone further development of this set-valued duality but sketch the opposite-variance bi-sheaf (48), associated with mutually adjoint contravariant functors. The pairings are proposed as the cross-Hom sets in a category $\mathcal E(A)$ combining $A^\wedge$ and $(B^\wedge)^o$ along A, with $\mathcal E(A)^o\simeq\mathcal E(B)$. Their compatibility is to be expressed through composition and associativity; full details are not supplied here. The later discussion on pp.\ 404--405 returns to this construction, as the marginal pointers anticipate (p.\ 398). Case~2, opened at the foot of p.\ 398, takes abelian-group coefficients and uses the abelianization results (12) and (14) of \S104. For $B=A^o$, (49) identifies $B^\wedge_{ab}$ with the category of additive cosheaves on $A^\wedge_{ab}$ valued in $(Ab)$. The corresponding pairing (50) commutes with small inductive limits in each argument, and its symmetry supplies the analogous description of $A^\wedge_{ab}$. The equivalence (51) reduces such pairings to their values on $A\times B$. Formula (52) gives the free abelian group on $\mathrm{Hom}(b,a)$ when the arguments are the abelianizations of representables; (53) extends this to $\underline Z^{(F*G)}$ for presheaves $F,G$. The insertion on p.\ 400 explicitly says that \emph{both members}, rather than only the left side, commute with the relevant inductive limits (pp.\ 399--400). Remark~1 abstracts this situation to two additive categories $\mathcal P,\mathcal Q$ with a pairing to $(Ab)$, cocontinuous in each variable, which identifies each category with the additive cosheaves on the other. The presheaf examples are abelian topoi: abelian categories with exact small filtered inductive limits and a small generating subcategory. The visible revisions add the filtering qualification and clarify that this terminology does not mean a category which is simultaneously an ordinary topos and an abelian category. Replacing abelian groups by modules over a commutative ring gives the corresponding $k$-linear examples. For a noncommutative ring, the two sides use right and left modules; the proposed pairing is the abelian-group pairing followed by quotienting the resulting bimodule by the differences $sx-xs$. Grothendieck explicitly says that he has not checked that this construction gives the desired duality. He does verify the ordinary one-object example through tensor product, evaluation at the regular module, and the associated tensor--Hom equivalences. Their comparison is an additive analogue of (39) when the target category has both kinds of small limits (pp.\ 400--402). Remark~2 singles out the existence of a small generating family of ultraprojective objects. For a module category, one ultraprojective generator suffices; ultraprojectivity is decomposed into projectivity and finite presentation. The proposed general structure theorem describes the resulting “algebraic” abelian topoi as categories of additive functors from $P^o$ to $(Ab)$, for a small additive category $P$. One can take $P$ essentially small and canonically equal to the full subcategory of ultraprojective objects. For $A^\wedge_{ab}$, this is the abelian Karoubi envelope of $A$, or alternatively one may use $\mathrm{Add}(A)$ for more explicit computation; the text refers back to \S\S93 and 99 (p.\ 402). The notation $P^{\&}$ and the inclusion $\varepsilon_P:P\to P^{\&}$ lead to the expected equivalence (**) between cocontinuous additive functors out of $P^{\&}$ and additive functors out of $P$. The suggested proof imitates the nonadditive representable-density argument. Its dual gives the projective-limit version and then (***), comparing sheaves and cosheaves on the dual algebraic abelian topoi. These statements are presented as the expected additive theory, with the proof still to be checked (p.\ 403). The next page calls an abelian topos \emph{reflexive} when it belongs to such a dual pair. Algebraicity appears sufficient, and the author wonders whether it is necessary, at least for the stronger equivalence with arbitrary additive target categories and suitable compatibility as the target varies. Remark~3 records the central limitation: for a general dual pair he does not see how to construct the comparison functor from the pairing alone. The same perplexity occurs in the nonabelian case, so the algebraic examples should not be read as an established general duality theorem (p.\ 404). Remark~4 returns to the auxiliary category $\mathcal E(A)$ from p.\ 398. Inside the functors from $B^\wedge$ to sets, it is represented by the union of the cocontinuous functors and the representable functors. Their intersection is the Karoubi envelope of $A$, equivalently the opposite of the ultraprojective objects of $B^\wedge$. The margin note describes the union by commutation with either type of small limits. These constructions are then said to extend to a small additive $P$ (pp.\ 404--405). Case~3 resumes the intended application of (39). With an additive target $\mathcal M$, initially under the extra pseudotopos assumptions on $\mathcal M$ and its opposite, (12) and (13) yield (54), relating $\mathcal M$-valued cosheaves on $A^\wedge_{ab}$ to sheaves on $B^\wedge_{ab}$. The earlier discussion suggests that those extra assumptions should be unnecessary. The required assertion (55) says that precomposition with abelianization gives the expected equivalence for any additive $\mathcal M$ with small inductive limits. Formula (56) reduces it to additive functors on $\mathrm{Add}(A)$, and (57) is the projective-limit dual. The four-arrow diagram (58) shows that the two left comparisons are equivalent conditions because the two right comparisons are already equivalences. Finally, identifying $A^\wedge_{ab}$ with $P^{\&}$ for $P=\mathrm{Add}(A)$ reduces the remaining task to (**). The author prefers to prove that general statement and ends by saying that the suggested proof should be checked during later scratchwork. The scan's handwritten reference in the insertion is (56), and its final wording is “proof given” followed by the correction from “there” to “here” (pp.\ 405--406).

\paragraph{\S106. Dual operations, a formulaire and explicit computation (pp.\ 407--419).}
The new subsection E), dated 14 and 15 August, organizes the preceding duality into a calculus for the operations $*$ and Hom. Formula (59) relates $\mathcal M$-valued cosheaves on $A^\wedge_{ab}$ to sheaves on the opposite presheaf topos. Formula (60) instead reverses the values by passing from $\mathcal M$ to $\mathcal M^o$ while retaining the same base. Combining them gives (61), a duality expressed entirely in terms of sheaves; the author prefers this presentation for the intended homology and cohomology applications. The enlargement of diagram (58), and its dual (62), compare functors on $\mathrm{Add}(A)$, its completion under small sums, the abelianized topoi, and their opposites. The text observes that the ten displayed occurrences represent nine distinct categories, using the canonical identification $\mathrm{Add}(B)^o\simeq\mathrm{Add}(B^o)=\mathrm{Add}(A)$. The dual diagram and its introducing sentence consistently use the handwritten coefficient $\mathcal N$, an important distinction from the earlier transcription's $\mathcal M$. Formulas (63)--(64) return to the abelian pairing and its interpretation as a cosheaf (pp.\ 407--408). Tensor structures on the two abelian topoi provide additional operations. Formula (65) is the internal Hom action of an abelian sheaf on an $\mathcal M$-valued sheaf. Through the duality, (66)--(67) express the corresponding covariant operation, denoted by a circled star. In the source this sign is repeatedly a typed star surrounded by a hand-drawn circle. The final-category example reduces these constructions to familiar Hom and tensor-product operations. Their commutation with small limits is not automatic: the text distinguishes the finitely presented projective situation from the weaker flatness detected by exactness of tensor product. The proposed converse from exactness to projectivity is explicitly rejected using flat versus free abelian groups (pp.\ 409--410). The next page retracts the initial feeling that the Hom side was already computationally understood. The data (68) are just the functors $F:A^o\to\mathcal M$, $L:A^o\to(Ab)$ and $L':B^o\to(Ab)$. Passing to the opposite coefficient category and interchanging $A,B$ produces (68') and the tautological identifications (69). Thus an adequate computation and formulaire for either operation should yield one for the other by duality; this is a program for the following pages, rather than a claim that the computation was already available (p.\ 411). The author first broadens the Hom calculation to a topos $\mathcal A$, and provisionally to suitable pseudotopoi. The coefficient category is additive and stable under small projective limits; the additional pseudotopos hypothesis on $\mathcal M$ is used when $\mathcal A$ is not a presheaf category. External Hom, with values in $\mathcal M$, is distinguished from internal Hom, with values in $\mathcal A_{\mathcal M}$, in (70). A sheaf $F$ is viewed as the projective-limit-compatible functor (71). The abelianization $X\mapsto\underline Z^{(X)}$, its componentwise form for a presheaf topos, the unit $\underline Z_{\mathcal A}=\underline Z^{(e)}$ in (72), and the sections functor (73)--(74) supply the ingredients for the main list (75) (pp.\ 412--413). Three entries of (75), circled on the page, are basic: a) evaluates external Hom on an abelianized representable; e) gives the relevant projective-limit compatibility in the first argument; and b) is the tensor--Hom relation. The remaining entries express their internal-Hom versions, the unit identities and the relation between internal Hom and sections. The text derives them through (76), the compatibility of abelianization with products, associativity of tensor product, and specialization at the final object. The comments explicitly leave aside coherence among the canonical isomorphisms and postpone change of topos to a later section connected with homology and cohomology invariants of maps in $(Cat)$ (pp.\ 413--414). For computation, a site $A$ presenting $\mathcal A$ allows $F$ and $L$ to be treated as the functors (77)--(79). When representables are sheaves, the canonical functor (78) is fully faithful; the text permits a change of site to reach that situation. The additive envelope and the extension $\varepsilon_{ab}:\mathrm{Add}(A)\to\mathcal A_{ab}$ in (80) turn the evaluation formula into the finite-product expression (81). The handwritten insertion and rewritten formula call the extension $\mathrm{Add}(F)^*$, although the subsequent expression (83) writes $\mathrm{Add}(F)$ without the star. A inductive limit presentation of $L$ by objects of $\mathrm{Add}(A)_{/L}$, compared with p.\ 403, and exactness e) then give (83). Its arrow under lim points \emph{left}: the displayed expression for $\mathrm{Hom}(L,F)$ is a projective limit, in contrast to the arXiv witness's inductive-limit arrow. The source also writes $x_i$ in the live tuple describing objects of the indexing category, then $a_i$ in the explanation; this inconsistency is preserved rather than silently repaired (pp.\ 415--416). The remark after (83) separates the construction from the provisional assumptions. The formula makes sense for an additive $\mathcal M$ with small projective limits, without a pseudotopos hypothesis on $\mathcal M$. This suggests an extension of the equivalence from (13), p.\ 388, when $\mathcal A$ is an actual topos. The proposed quasi-inverse still has to be checked on the abelianized representables. That check is supplied in the presheaf case by full faithfulness of (80), which gives a final object in the relevant indexing category. The author then regards the immediate computational program as complete and turns to its dual, while the proposed general extension remains qualified by the stated verification (p.\ 416). Dualizing gives the list (84), now under the assumption $\mathcal A=A^\wedge$ and with small inductive limits in $\mathcal M$. Its external star and internal circled-star formulas parallel (75), but the source retains a plain star in e'). The asymmetry of the a') formulas comes from the possible absence of binary products or sums in $A$. Formula (85) spells out the Hom calculation using the category over $a\times b$; the dual symbol $a\vee b$ initially denotes an object of $(B^\wedge)^o$ and agrees with the sum in $A$ when that sum exists. Formula (86) describes the dual indexing category concretely by triples $(x,u,v)$ with maps $a\to x$ and $b\to x$ (pp.\ 417--418). Remark~1 treats additive $A$, where finite sum and product agree. The antiequivalence (88) gives the comparison (87) between internal Hom and circled-star for $L\in\mathrm{Add}(A)$, compatible with tensor products. A faint margin note proposes extension to $\mathrm{Kar}\,\mathrm{Add}$; it is not a proof of that extension. Remark~2 begins the comparison with two more elementary operations by introducing componentwise tensor product (89); its dual is taken up on the next page (p.\ 418). The closing page gives the dual componentwise operation (90), taking external Hom at each object and requiring small projective limits in the additive coefficient category. Both elementary operations extend to arbitrary topoi, with abelian cosheaves replacing the opposite presheaf category in the dual construction. The author does not develop them here, but points to the role of tensor product and its derived operation for ringed topoi and anticipates a parallel role for (90). The final paragraph explains the reuse of the Hom notation: its coefficient argument lies in different categories in (65) and (90), while in the final-category case the operations agree canonically. The same reasoning applies to tensor product and the circled-star operation. This closes the notation review before the dated transition, 17 August, to ground-ring extension in \S107 (p.\ 419).

\paragraph{\S107. Ground-ring extension, additive envelopes and resolutions (pp.\ 419--425).}
Subsection F), dated 17 August, extends the preceding constructions from $\underline Z$ to a commutative ring $k$. Formula (91) specifies a $k$-additive category by an action of $k$ on the identity functor of an additive category. The compatible functors form the full subcategory (92) of the additive functors in (93). Lower and upper exclamation indices retain their respective direct- and projective-limit meanings; bilinearity means $k$-additivity in each variable. The most intrinsic possible ground ring for a given additive category is the endomorphism ring of its identity functor, and other scalar structures are obtained by restriction along a ring map. Formula (94) then defines $A^\wedge_k$ as the category of presheaves of $k$-modules on $A$ (pp.\ 419--420). A homomorphism $k\to k'$ induces the componentwise scalar-extension functor (95), $F\mapsto F\otimes_k k'$. The text interprets $A^\wedge_{k'}$ as a ground-ring extension of $A^\wedge_k$ through a universal problem for cocontinuous $k$-linear functors into $k'$-additive targets. The $k$-abelianization (96), $X\mapsto k^{(X)}$, is the composition of ordinary abelianization with extension from $\underline Z$ to $k$. Precomposition gives (97), identifying cocontinuous $k$-linear functors on $A^\wedge_k$ with the corresponding functors on $A^\wedge$, and hence with arbitrary functors from $A$. Formula (98) is the projective-limit dual. Replacing $A$ by $B=A^o$ and assuming both kinds of small limits in the coefficient category produces the duality (99), interpreted through $\mathcal M$-valued sheaves on $B$ (p.\ 421). The two descriptions induce the external star and internal circled-star operations (100), and the external and internal Hom operations (101). Their rules are inherited from (84) and (75), with $k$ subscripts. A typed insertion supplies the limit hypothesis; its coefficient-category position is left blank, while the handwritten replacement specifies inductive limits for the star side and projective limits for the Hom side. Explicit computation requires the $k$-additive envelope $\mathrm{Add}_k(A)$ in (102), consisting of finite sums of the objects $k^{(a)}$, and its extension $\mathrm{Addinf}_k(A)$ in (103), allowing arbitrary small sums. The latter solves the universal problem of mapping $A$ into an infinitely additive $k$-linear category. The source briefly writes $A^\wedge$ without the lower $k$ in the prose description of the first envelope; that omission is preserved in the edition rather than supplied from the surrounding formulas (p.\ 422). Adjoining direct factors gives the projective objects when one starts from $\mathrm{Addinf}_k(A)$, and the ultraprojective objects when one starts from $\mathrm{Add}_k(A)$. Ultraprojectivity here means projectivity together with finite presentation. Formula (104) identifies these categories with the respective Karoubi envelopes; the printed operator for the latter subcategory is $\mathrm{ULProj}$. Thus the five $k$-additive categories are characterized by different combinations of finite or arbitrary sums, images of projectors, and small inductive limits. Adding $A$ and $A^\wedge$ gives the seven-object diagram (105). The handwritten qualification says that the functors among the five middle categories are fully faithful (p.\ 423). Restriction of cocontinuous $\mathcal M$-valued functors gives the seven-term diagram (106). Its two middle horizontal comparisons are drawn as three-stroke equivalence signs without arrowheads; the other five comparisons are arrows. The suffixes ``add'', ``addinf'' and ``kar'' encode the corresponding preservation requirements. For $k=\underline Z$, part of this diagram recovers (58); its projective-limit dual contains (62). When $\mathcal M$ has both kinds of small limits, the last terms coincide, and the author describes the combined system as fourteen mutually equivalent categories, giving alternative interpretations of an $\mathcal M$-valued copresheaf on $A$ (p.\ 424). The closing discussion separates the purposes of the five additive categories. The ambient presheaf category supports general $k$-linear constructions; its projectives support resolutions. The equivalence (*) between the derived category bounded above and the homotopy category of bounded-above projective complexes extends in (107) to complexes in $\mathrm{Addinf}_k(A)$. Direct sums of the $k^{(a)}$ are convenient because their morphisms can be described by matrices with entries in free modules on the sets $\mathrm{Hom}(a,b)$. A cointegrator with coefficients in $k$ is a projective resolution of the constant presheaf $k_A$, viewed up to the stated unique isomorphism in these homotopy categories. Complexes with components in $\mathrm{Addinf}_k(A)$ are called quasi-special; those with components already in $\mathrm{Add}_k(A)$ are special. Formula (108) identifies the homotopy categories formed from $\mathrm{Add}_k(A)$ and its Karoubi envelope $\mathrm{ULProj}(A^\wedge_k)$. The claim that their canonical functor into (107) is fully faithful remains explicitly tentative, subject at least to suitable coherence conditions; it is not presented as an unconditional conclusion (pp.\ 424--425).

\paragraph{\S108. Homology, cohomology\newline and the limits of the projective construction (pp.\ 426--432).}
The entry dated 22 August begins with a personal interruption: the death of the author's granddaughter Ella and the earlier death of Ronnie Brown's son Gabriel. His response to Brown's letter frames mathematical correspondence as a possible human relationship as well as an exchange of ideas. Subsection G), on homology and cohomology in the absolute case, then resumes the review. Its organizing aim is a precise duality between the two theories, requiring coefficients in general abelian categories with the relevant small direct or projective limits. Passing to the opposite coefficient category and to $B=A^o$ should interchange homology and cohomology (p.\ 426). Cohomology supplies the familiar starting point: sheaf cohomology is expressed by Ext from the constant integral sheaf, or from the structure sheaf on a ringed topos. The author fixes a commutative ground ring $k$ for the intended duality and notes, without pursuing it, the possibility of a locally constant sheaf of ground rings. He emphasizes that the total derived object carries more information than its individual Ext groups. Pairings (*) and (**) recall external and internal derived Hom, and a projective resolution of the first argument computes them in the presheaf situation. Taking the first argument to be the structure sheaf gives derived sections in (i); (ii) relates external derived Hom to derived sections of internal derived Hom. These are recalled first in the familiar ringed-topos setting, before the more general coefficient problem is attempted (pp.\ 427--428). The ordinary pairing (109) takes a $k$-module presheaf and an $\mathcal M$-valued presheaf to an object of $\mathcal M$. The proposed derived pairings (110) and (111) have bounded-above first argument and bounded-below second argument; the second pairing takes values in the derived category of $\mathcal M$-valued presheaves. Abelianity of $\mathcal M$ is needed to speak of these derived categories. Formula (112) names the cohomology objects as external and internal Ext. The source's handwritten arguments in its second expression occur as $F_\bullet,L^\bullet$, and this literal order is retained in the edition. The distinction between external and internal refers to the value category, while the two arguments already belong to different presheaf categories (p.\ 428). Because enough injectives in the coefficient category have not been established, (113) attempts to construct both pairings using a projective resolution of the first argument. The resulting double-Hom simple complexes are defined up to homotopy even before the abelian hypothesis. To descend with respect to quasi-isomorphisms in the second argument, Hom from each fixed projective must be exact. For external Hom this follows first for a representable $k^{(a)}$ by evaluation, then for small sums by the hypothesis that small products of epimorphisms are epimorphisms, and finally for their direct factors. Thus the external pairing (110) is obtained under the stated exactness assumption. The analogous internal calculation involves $F(a\times b)$ and immediately exposes a difficulty when $A$ is not stable under binary products (p.\ 429). The next page explicitly retracts the proposed projective computation of internal derived Hom. Projectivity of a sheaf of modules is not generally preserved by localization, in contrast to the injective behavior recalled there. Localizing $k^{(a)}$ over $b$ leads to the object represented by $a\times b$ in the slice presheaf topos; the displayed Ext calculation need not vanish in positive degree. Such vanishing would impose strong acyclicity conditions on the connected components of that product. Consequently (111) has not been constructed by the proposed argument. The possibility of using enough injectives is postponed. External derived Hom is nevertheless sufficient for (114): derived sections are computed by Hom from a projective resolution (115) of $k_A$. Such a resolution is the cointegrator, and (116) is the resulting functor on bounded-below complexes of $\mathcal M$-valued presheaves. The canceled underline beneath Hom in (114) visibly reinforces this distinction between external and internal constructions (p.\ 430). The author acknowledges that the expected comparison (ii) has therefore not been obtained in this generality, then dualizes the construction that survives. The coefficient category $\mathcal N$ is abelian, has small inductive limits, and satisfies the dual exactness condition on small sums of monomorphisms. Formula (117) is the derived external star pairing; (118) computes it by resolving its second argument in the opposite presheaf category. The composite equivalence (119) and the tautological isomorphisms (120), (120') express the interchange with external derived Hom under $\mathcal M=\mathcal N^o$ and $B=A^o$. The two occurrences of $A^\wedge$ in (119) have no coefficient subscript in the scan. Integration is then described, with the inserted qualification ``at least morally'', as total left derivation of inductive limit; (121) and the number corrected by hand to (121') suggest the dual derived-limit notation (p.\ 431). The closing page seeks notation that retains the geometric associations of homology and cohomology while displaying their symmetry. Formula (123) introduces $\mathrm{RH}^{\bullet}$ and $\mathrm{LH}_{\bullet}$, and (124) states their interchange on passing to the dual base and opposite coefficients. The displayed targets of (123) are literally $\mathcal M$ and $\mathcal N$; the edition does not silently supply derived-category symbols. The final remarks delimit the achievement. The symmetry uses fairly general abelian coefficient categories satisfying the specified exactness properties, but $\mathrm{RH}^{\bullet}$ has not yet been shown to be the total right derived functor of sections or projective limit. The notation (121), (121') is consequently called misleading. Establishing enough injectives and comparing the injective and projective definitions should also recover a satisfactory internal derived Hom formalism. This unresolved requirement leads directly into the next day's investigation in \S109 (p.\ 432).

\paragraph{\S109. Projectives, injectives and a return to derived Hom\newline (pp.\ 433--442).}
The entry dated 23 August answers the question left open in \S108 by working first with a small $k$-additive category $P$. Formula (125) defines $P^{\&}_{\mathcal N}$ as the category of $k$-additive functors from $P^o$ to $\mathcal N$; its small limits and abelian structure are computed componentwise when those structures exist in $\mathcal N$. Proposition~4 says that small inductive limits and enough projectives in $\mathcal N$ give enough projectives in $P^{\&}_{\mathcal N}$. If those projectives are also $k$-flat, a projective coefficient presheaf makes the relevant star pairing exact in its other argument. The assumptions are necessary already when $P$ is the additive category generated by the final category. The special case $P=\mathrm{Add}_k(A)$ recovers the presheaf categories needed in the preceding section (pp.\ 433--434). The proof uses the external product $L\otimes x$ of (126), computed componentwise. The three Hom expressions in (127) give its tensor--Hom adjunction; the comparison involving the first expression requires small projective limits in $\mathcal N$, while the direct comparison of the other two does not. It follows that the tensor product of projective arguments is projective. Formula (127') specializes to a representable and gives the evaluation adjunction. Choosing projective covers of the values $F(a)$ then produces a projective epimorphism onto $F$ from the sum of the corresponding representables tensored with those covers. Compatibility (129) proves the supplementary flatness assertion. Its necessity is illustrated by $k=\underline Z$ and coefficients in vector spaces over $\mathbf F_p$: every coefficient object is projective, while only the zero object is flat over $\underline Z$. The dual statement gives enough injectives and the corresponding $k$-coflatness assertion. Several coefficient-category positions are literally blank in the scan and remain blank in the transcription (pp.\ 434--435). Corollary~1 returns to $A$ and lists four conditions on $\mathcal M$: abelianity with small projective limits, enough injectives, $k$-coflatness of its injectives, and preservation of epimorphisms by small products. Under these hypotheses, external derived Hom in (130) can be computed by a projective resolution of its first argument, an injective resolution of its second argument, or both; (131) expresses the compatible constructions. The dual statements (130'), (131') concern derived star and projective coefficient resolutions. This settles the availability of the resolutions required for the external pairing and opens the way to the internal problem (p.\ 436). The internal operations (132), (132') still require care. Projective first arguments alone do not generally suffice, as \S108 already showed. Resolving the second argument by injectives does work if tensoring with a representable is exact in the relevant sense. For $P=\mathrm{Add}_k(A)$, componentwise tensor product makes representables, and their finite sums, $k$-flat. Corollary~2 therefore constructs internal derived Hom (133) from an injective resolution of its second argument. Formula (134) gives the desired comparison: external derived Hom is obtained by applying derived cohomology to internal derived Hom. The circled-star dual (135) uses a projective resolution of its coefficient argument, and (134') compares it with derived homology. The arrow into the target category is absent from the displayed (135) in the source; the edition preserves that omission. These results supply the internal formalism that the earlier projective construction had failed to provide (pp.\ 437--438). Remark~1 broadens the setting again. For a general small additive $P$, an internal tensor structure would come from a diagonal $P\to P\otimes_k P$, together with the appropriate coalgebra structure. An ordinary category supplies such a diagonal before linearization. The author proposes that a coalgebra object in the $2$-category of small $k$-additive categories might give a finer model of homotopy types than the De Rham constructions with divided powers considered earlier, referring back to \S94, p.\ 339. Associativity and commutativity should be specified by coherent isomorphisms, rather than imposed as identities; they are additional data with coherence conditions. The references to Saavedra's thesis and to work of Mac Lane, Bénabou and Sinh indicate hoped-for sources of this formalism. The text presents this as an avenue to investigate and explicitly hopes that reversing arrows will yield the parallel theory (pp.\ 438--439). The next page develops some consequences of that proposal. Extension from $k$ to $k'$ should preserve the objects of $P$ and tensor its Hom modules with $k'$, with the stated $2$-categorical universal property. Units $k_P$ and $k_Q$ correspond to augmentation and coaugmentation data. Derived Hom from the former and derived star with the latter should recover cohomology and homology, with constant coefficients and cup and cap operations interpreted through the tensor structures. Coalgebra objects could also be considered at the derived-category level. The author does not expect this alone to recover a general fundamental group or homotopy type: outside the $1$-connected case he anticipates a groupoid action on the universal-cover data. The following page asks whether a category of homotopy types over $k$, compatible with base change, might be approached this way without first constructing the full proposed models (pp.\ 440--441). Remark~2 sharpens the hypotheses of Corollary~1. The product exactness condition is used for the projective-resolution construction, while enough injectives and their coflatness support the injective construction; abelianity with projective limits supplies the ambient categories. When all four conditions hold, the two constructions agree. Cointegrators remain the author's preferred method for external cohomology, but internal derived Hom uses injective coefficient resolutions. In particular, the first three conditions suffice for the internal construction and the comparison (134); the product condition is not required when one works throughout with injectives. The source calls the injective complex $K^\bullet$ in this discussion, and that literal letter is retained (p.\ 441). Remark~3 closes the section with the general diagonal case. Exactness of $L\mapsto L\otimes a$ follows if each $\mathrm{Hom}_P(b,a)$ is $k$-flat; equivalently, the projectives in $P^{\&}$ must be flat for the external tensor action. In the case $P=\mathrm{Add}_k(A)$, the relevant Hom modules are finite sums of free modules on the sets of morphisms in $A$, hence projective. The author would prefer projectivity in the general theory, but flatness suffices for the asserted injectivity of external and internal Hom when the first argument is projective and the second injective. Finally, for a $k$-additive functor $u:\mathcal M\to\mathcal M'$ commuting with small projective limits, (136) compares $Ru$ applied to derived Hom with derived Hom after the componentwise derived functor $Ru^P$. If $u$ is exact, the derivation of $u$ is unnecessary. The analogous internal and dual formulas are explicitly omitted. This completes \S109 and the cote, before the next day's \S110 returns to schematization (p.\ 442).
